\documentclass[conference]{IEEEtran}

% Packages
\usepackage{graphicx}
\usepackage{hyperref}
\usepackage{amsmath}
\usepackage{amssymb}
\usepackage{booktabs}

% Hyperref setup
\hypersetup{
    colorlinks=true,
    linkcolor=blue,
    urlcolor=blue,
    citecolor=blue
}

\begin{document}

\title{Informe del Artículo: Positional versus Symbolic Attention Heads}

\author{
    \textbf{Autores originales:} Felipe Urrutia, Juan Jos\'e Alegr\'ia,\\
    Cinthia Sanchez Macias, Jorge Salas, Cristian B. Calderon,\\
    Cristobal Rojas\\
    \textbf{Afiliaci\'on:} CENIA \& Faculty of Mathematics UC, Santiago, Chile\\
    \textbf{arXiv:} 2605.31558 (NeurIPS 2026)
}

\maketitle

\begin{abstract}
Este art\'iculo estudia c\'omo los \textbf{attention heads} en transformers desarrollan comportamientos \textbf{posicionales} versus \textbf{simb\'olicos} durante el entrenamiento, y c\'omo estos comportamientos afectan la capacidad de generalizaci\'on a secuencias m\'as largas. Se introducen tareas estructuralmente equivalentes que requieren perfiles de atenci\'on distintos, se define una descomposici\'on mecanicista en funciones b\'asicas (Index, Retrieval, Reflexive), y se presenta un an\'alisis te\'orico de c\'omo RoPE puede implementar estas funciones con interpretaci\'on geom\'etrica, estableciendo cotas de discrepancia que explican cuantitativamente la degradaci\'on con la longitud de secuencia.
\end{abstract}

\section{Motivaci\'on y Contexto}

Los transformers modernos (GPT, Claude, etc.) usan \textbf{Rotary Positional Encoding (RoPE)} para codificar informaci\'on posicional. Sin embargo, no se entend\'ia bien c\'omo los attention heads aprenden a usar la informaci\'on posicional (atender a posiciones espec\'ificas) versus informaci\'on simb\'olica (atender a tokens espec\'ificos sin importar su posici\'on).

Trabajo previo de Urrutia et al. (2025) introdujo \textbf{m\'etricas} para clasificar heads como posicionales o simb\'olicas, y observaron que en modelos reales las capas tempranas tienden a ser posicionales y las profundas simb\'olicas. Este art\'iculo \textbf{extiende ese an\'alisis al estudio de la din\'amica de aprendizaje} y a la \textbf{generalizaci\'on a largo contexto}.

\section{Tareas Propuestas}

\subsection{Number Task (Tarea Num\'erica)}

\begin{itemize}
    \item \textbf{Entrada:} Una secuencia de letras seguida de n\'umeros.
    \item \textbf{Mecanismo:} Cada n\'umero indica cu\'antas posiciones retroceder en la secuencia.
    \item \textbf{Ejemplo (4-hop):} \texttt{C B A ... 1 4 1 3 2 1 1} $\rightarrow$ hop1: retrocede 1 $\rightarrow$ hop2: retrocede 4 $\rightarrow$ hop3: retrocede 1 $\rightarrow$ hop4: retrocede 3 $\rightarrow$ respuesta: letra en posici\'on final.
    \item \textbf{Tipo de razonamiento:} \textbf{Posicional} -- requiere seguir posiciones relativas.
    \item \textbf{Vocabulario:} 120 letras (a--z + bigramas aa--dp) + 16 enteros $\{1,\dots,16\}$ = 136 tokens.
\end{itemize}

\subsection{Letter Task (Tarea de Letras)}

\begin{itemize}
    \item \textbf{Entrada:} Una secuencia de pares letra-letra seguida de pares letra-n\'umero.
    \item \textbf{Mecanismo:} Cada par $(X,Y)$ indica ``busca el par cuya primera letra es $Y$''.
    \item \textbf{Ejemplo (4-hop):} \texttt{(I,G) (G,D) (D,E) (E,C) (H,F) ... (C,2)} $\rightarrow$ hop1: $G$ busca a $(I,G)$ $\rightarrow$ hop2: $D$ busca a $(G,D)$ $\rightarrow$ \dots $\rightarrow$ respuesta: $(C,2)$.
    \item \textbf{Tipo de razonamiento:} \textbf{Simb\'olico} -- requiere seguir asociaciones entre s\'imbolos.
    \item \textbf{Vocabulario:} 8 letras $\{a,\dots,h\}$ + 8 enteros $\{1,\dots,8\}$ = 64 pares $\times$ 2 tipos = 128 tokens.
\end{itemize}

Ambas tareas son \textbf{estructuralmente equivalentes} (misma estructura multi-hop), pero demandan \textbf{mecanismos distintos}.

\section{Configuraci\'on Experimental}

\subsection{Arquitectura del Modelo}

\begin{itemize}
    \item \textbf{Modelo:} GPT-J (decoder-only Transformer)
    \item \textbf{Capas:} 12
    \item \textbf{Attention heads:} 1 por capa (single-head, para aislar mecanismos)
    \item \textbf{Dimensi\'on oculta:} 128
    \item \textbf{Posicional encoding:} RoPE
    \item \textbf{Activaci\'on:} GELU
    \item \textbf{Dropout:} 0.1
    \item \textbf{Objetivo:} Last-token prediction (solo se optimiza la p\'erdida en el \'ultimo token)
\end{itemize}

\subsection{Dataset}

\begin{itemize}
    \item \textbf{480,000 secuencias} por tarea
    \item Cada secuencia: 17 tokens (8 target window + 8 pre-target window + 1 query)
    \item Balanceado: misma proporci\'on de casos 1-hop, 2-hop, 3-hop, 4-hop
    \item Split: 90\% train / 0.5\% validation / 9.5\% test
\end{itemize}

\subsection{Entrenamiento}

\begin{itemize}
    \item Entrenamiento con 3 NVIDIA RTX A5000 (24GB cada una)
    \item Tiempo total: $\sim$55 minutos por modelo
    \item M\'etricas de positional/symbolic score: 10h GPU + 17h CPU
\end{itemize}

\section{M\'etricas: Positional y Symbolic Scores}

Basado en Urrutia et al. (2025), se definen dos m\'etricas para cuantificar el comportamiento de un attention head.

\subsection{Definici\'on Formal}

Un head act\'ua \textbf{posicionalmente} si su distribuci\'on de atenci\'on es \textbf{invariante} ante permutaciones de los key vectors. Es decir, atiende a posiciones espec\'ificas sin importar qu\'e token est\'e ah\'i.

Un head act\'ua \textbf{simb\'olicamente} si su distribuci\'on de atenci\'on es \textbf{equivariante} ante permutaciones. Es decir, atiende a tokens espec\'ificos sin importar su posici\'on.

\subsection{Scores}

Se calculan mediante:
\begin{itemize}
    \item \textbf{Permutaciones por pares} (swaps) entre posiciones $i$, $j$
    \item \textbf{Peso de cada permutaci\'on} basado en la masa de atenci\'on movida (softmax con temperatura)
    \item \textbf{Positional Score:} Similitud coseno entre atenci\'on original y atenci\'on tras permutar
    \item \textbf{Symbolic Score:} Similitud coseno entre atenci\'on original y atenci\'on ``can\'onica'' tras permutar
\end{itemize}

\section{Resultados Principales}

\subsection{Din\'amica de Aprendizaje (Figura 2)}

\textbf{Number Task:}
\begin{itemize}
    \item El accuracy emerge \textbf{progresivamente}: primero se aprende 1-hop, luego 2-hop, etc.
    \item Se requieren \textbf{heads posicionales} para el Indexing + \textbf{heads simb\'olicas} para propagaci\'on.
    \item Las heads posicionales aparecen en capas medias (L3--L6).
    \item Las heads simb\'olicas aparecen en capas profundas (L9--L11).
\end{itemize}

\textbf{Letter Task:}
\begin{itemize}
    \item El accuracy emerge \textbf{simult\'aneamente} para todos los hops.
    \item Solo se requieren \textbf{heads simb\'olicas} (Retrieval).
    \item La \'ultima head simb\'olica en L11 es necesaria para todos los hops.
    \item Las heads posicionales no son necesarias.
\end{itemize}

\subsection{Pureza de Heads}

La \textbf{pureza} (heads que son claramente posicionales O simb\'olicas, no mixtas) es necesaria para alcanzar m\'axima precisi\'on. Esto valida las m\'etricas como herramienta de interpretabilidad.

\subsection{Funciones Mecanicistas Identificadas}

Tres funciones b\'asicas implementadas por los heads puros:

\begin{enumerate}
    \item \textbf{Selective Index (Index):} Dada una posici\'on con un n\'umero $n$, copia el token $n$ posiciones atr\'as, pero solo si es una letra. $\rightarrow$ \textbf{Posicional}

    \item \textbf{Retrieval:} Dado un par $(X,Y)$, busca el par anterior cuya primera letra sea $Y$ y lo copia. $\rightarrow$ \textbf{Simb\'olico}

    \item \textbf{Reflexive:} Copia el token actual a la siguiente capa sin cambios. $\rightarrow$ Puede ser posicional o simb\'olico (en la pr\'actica, simb\'olico).
\end{enumerate}

\subsection{Descomposici\'on Mecanicista}

Para una entrada con $h$ hops y un modelo con $\ell \ge h$ capas:

\begin{itemize}
    \item \textbf{Number Task:} $f_{\text{Num}}(\sigma) = \text{Reflexive}^{(\ell-h)}(\text{Index}^{h}(\sigma))[n]$
    \item \textbf{Letter Task:} $f_{\text{Let}}(\sigma) = \text{Reflexive}^{(\ell-h)}(\text{Retrieval}^{h}(\sigma))[n]$
\end{itemize}

\section{An\'alisis Te\'orico: Mecanismos Geom\'etricos en RoPE}

\subsection{Teorema 3.1 (Factibilidad)}

Existe un Transformer $T_{\text{Index}}$ con atenci\'on RoPE 2D que implementa Index, y un $T_{\text{Retrieval}}$ que implementa Retrieval, ambos con una sola cabeza de atenci\'on y sin funciones de activaci\'on.

\textbf{Mecanismo para Index (Posicional):}
\begin{itemize}
    \item Key vectors: todos alineados (iguales a un vector $v$)
    \item Query vectors: codifican la posici\'on a atender mediante rotaci\'on RoPE
    \item La atenci\'on m\'axima ocurre exactamente en la posici\'on deseada $(n - \sigma_n)$
    \item Los value vectors solo copian si el token es letra (mapeo diagonal)
\end{itemize}

\textbf{Mecanismo para Retrieval (Simb\'olico):}
\begin{itemize}
    \item Key vectors: codifican la letra izquierda del par mediante rotaci\'on con \'angulo $\omega$
    \item Query vectors: codifican la letra derecha del par consultante
    \item La atenci\'on m\'axima ocurre donde la letra izquierda del key coincide con la letra derecha del query
    \item La posici\'on exacta no importa, solo la identidad del s\'imbolo
\end{itemize}

\subsection{Remark 3.1}

Index \textbf{no puede} implementarse con un head puramente simb\'olico.
Retrieval \textbf{no puede} implementarse con un head puramente posicional.

Esto demuestra la \textbf{separaci\'on fundamental} entre ambos tipos de mecanismos.

\section{An\'alisis de Robustez: Discrepancia}

\subsection{Definici\'on de Discrepancia (Definici\'on 3.1)}

La \textbf{discrepancia $\Delta$} de un attention head es la diferencia m\'axima entre el logit m\'as grande y el segundo m\'as grande para entradas de cierta longitud.

Intuitivamente: mide qu\'e tan distinguible es el token correcto del resto. Una discrepancia de 0 significa que el modelo no puede distinguir el token correcto.

\subsection{Teorema 3.2 (Cotas de Discrepancia)}

\textbf{Para Index (posicional):}
\[
\Delta_{\text{Index}}(n) \leq \min\left\{1 - \cos(\theta), \frac{\pi^2}{2n^2}\right\}
\]

\textbf{Para Retrieval (simb\'olico):}
\[
\Delta_{\text{Retrieval}}(n) \leq 1 - \cos(\omega - n\theta)
\]

\textbf{Interpretaci\'on:}
\begin{itemize}
    \item Index: la discrepancia decae como $\mathcal{O}(1/n^2)$ -- se degrada r\'apidamente con la longitud.
    \item Retrieval: la discrepancia se mantiene \textbf{no nula} incluso para $n$ grandes si $\theta$ es peque\~no.
    \item Si $\theta = 0$ (NoPE), $\Delta_{\text{Retrieval}}$ se mantiene positiva para cualquier longitud.
\end{itemize}

Esto demuestra cuantitativamente que \textbf{los mecanismos simb\'olicos generalizan mejor} a secuencias largas.

\subsection{Validaci\'on Experimental (Figura 5)}

Los resultados confirman las predicciones te\'oricas:

\textbf{Modelo entrenado GPT-J:}
\begin{itemize}
    \item Letter task: $>90\%$ de precisi\'on hasta 850 tokens ($53\times$ la longitud original)
    \item Number task: la precisi\'on cae r\'apidamente despu\'es de la longitud de entrenamiento
\end{itemize}

\textbf{Modelos frontera (GPT 5.4, GPT 5.5, Claude Sonnet 3.7):}
\begin{itemize}
    \item Letter task: precisi\'on $>0.88$ a longitud 32, $>0.65$ a longitud 100
    \item Number task: precisi\'on $<0.5$ a longitud 32, $<0.1$ a longitud 100
    \item Consistente con las predicciones del Teorema 3.2
\end{itemize}

\section{Conclusiones y Contribuciones}

\subsection{Contribuciones Principales}

\begin{enumerate}
    \item \textbf{Tareas estructuralmente equivalentes} que requieren perfiles de atenci\'on distintos (posicional vs.~simb\'olico).
    \item \textbf{Relaci\'on entre pureza de heads y precisi\'on:} La m\'axima precisi\'on requiere heads puros (claramente posicionales o simb\'olicos).
    \item \textbf{Descomposici\'on mecanicista} de ambas tareas en funciones b\'asicas (Index, Retrieval, Reflexive).
    \item \textbf{Construcciones te\'oricas} que muestran c\'omo RoPE puede implementar estas funciones con interpretaci\'on geom\'etrica.
    \item \textbf{Noci\'on de discrepancia} para cuantificar la degradaci\'on con la longitud.
    \item \textbf{Separaci\'on cuantitativa} entre mecanismos posicionales y simb\'olicos en t\'erminos de generalizaci\'on.
\end{enumerate}

\subsection{Implicaciones}

\begin{itemize}
    \item Los mecanismos simb\'olicos son \textbf{intr\'insecemente m\'as robustos} para longitudes largas.
    \item Esto sugiere que entrenar modelos para que desarrollen estrategias simb\'olicas (vs.~posicionales) puede mejorar la generalizaci\'on.
    \item Las m\'etricas positional/symbolic scores pueden servir como \textbf{herramienta de monitoreo} durante el entrenamiento de LLMs.
\end{itemize}

\subsection{Limitaciones}

\begin{itemize}
    \item Configuraci\'on controlada con single-head por capa (no multi-head como modelos reales).
    \item Vocabulario finito limita pruebas de extrapolaci\'on.
    \item Solo se estudiaron tareas multi-hop simples.
\end{itemize}

\section*{Referencias Clave}

\begin{enumerate}
    \item Urrutia et al. (2025) --- ``Decoupling positional and symbolic attention behavior in transformers'' (arXiv:2511.11579) --- Introduce las m\'etricas positional/symbolic scores.
    \item Barbero et al. (2024) --- ``Round and round we go! What makes rotary positional encodings useful?'' --- Observa que heads simb\'olicas usan frecuencias bajas de RoPE.
    \item Su et al. (2024) --- ``RoFormer: Enhanced Transformer with Rotary Position Embedding'' --- Paper original de RoPE.
    \item Pasten et al. (2025) --- ``Continuity and isolation lead to doubts or dilemmas in large language models'' --- Sobre la dispersi\'on de atenci\'on con softmax.
\end{enumerate}

\end{document}
